% ===================================================================== PREAMBULE COMMUN — Carnet d'entraînement Fiche 9
\documentclass[11pt,a4paper]{article}
\usepackage[utf8]{inputenc}
\usepackage[T1]{fontenc}
\usepackage{lmodern}
\usepackage[french]{babel}
\usepackage{amsmath,amssymb,mathtools}
\usepackage{icomma}
\usepackage{siunitx}
\sisetup{output-decimal-marker={,},per-mode=symbol}
\DeclareSIUnit{\molar}{M}
\DeclareSIUnit{\Molar}{mol\,L^{-1}}
\usepackage[version=4]{mhchem}
\usepackage{xcolor}
\usepackage{colortbl}
\usepackage{tikz}
\usepackage{pgfplots}
\usepackage[most]{tcolorbox}
\usepackage{enumitem}
\usepackage{array}
\usepackage{booktabs}
\usepackage{multicol}
\usepackage{fancyhdr}
\usepackage[hidelinks]{hyperref}
\usetikzlibrary{arrows.meta,positioning,calc,decorations.pathmorphing,
  decorations.markings,angles,quotes,patterns,backgrounds,
  shapes.geometric,shapes.misc,fadings,intersections,fit}
\pgfplotsset{compat=1.18}
\usepackage[a4paper,margin=2.2cm,headheight=26pt,headsep=10pt]{geometry}

% ---------- Palette ----------
\definecolor{primary}{RGB}{150,98,162}
\definecolor{primarydark}{RGB}{92,52,110}
\definecolor{defcol}{RGB}{0,114,178}
\definecolor{thmcol}{RGB}{0,140,120}
\definecolor{formulacol}{RGB}{198,93,9}
\definecolor{remarkcol}{RGB}{120,94,200}
\definecolor{attncol}{RGB}{200,40,55}
\definecolor{excol}{RGB}{0,135,90}
\definecolor{methcol}{RGB}{90,90,100}
\definecolor{cationcol}{RGB}{0,90,190}
\definecolor{anioncol}{RGB}{200,60,55}
\definecolor{cristalcol}{RGB}{110,105,120}
\definecolor{eaucol}{RGB}{120,180,220}
\definecolor{solcol}{RGB}{0,135,90}
\definecolor{kpscol}{RGB}{198,93,9}
\definecolor{qcol}{RGB}{120,94,200}
\definecolor{precipcol}{RGB}{110,70,40}
\definecolor{exercol}{RGB}{150,98,162}
\definecolor{corrcol}{RGB}{0,120,100}

% ---------- Macros ----------
\newcommand{\Kps}{K_{\mathrm{ps}}}
\newcommand{\Ce}[1]{\left[\,\ce{#1}\,\right]}

% ---------- Style des boîtes ----------
\tcbset{boxbase/.style={enhanced,breakable,boxrule=0.9pt,arc=2.5pt,
   left=3mm,right=3mm,top=2.3mm,bottom=2.3mm,
   fonttitle=\bfseries,coltitle=white,toptitle=0.6mm,bottomtitle=0.6mm}}

\newtcolorbox[auto counter]{definition}[1][]{%
  boxbase,colback=defcol!5,colframe=defcol,
  title={\textsc{Définition}~\thetcbcounter\ \ifx\relax#1\relax\relax\else\textmd{--- \itshape #1}\fi}}
\newtcolorbox[auto counter]{theoreme}[1][]{%
  boxbase,colback=thmcol!6,colframe=thmcol,
  title={\textsc{Loi / Propriété}~\thetcbcounter\ \ifx\relax#1\relax\relax\else\textmd{--- \itshape #1}\fi}}
\newtcolorbox[auto counter]{formule}[1][]{%
  boxbase,colback=formulacol!6,colframe=formulacol,
  title={\textsc{Formule}~\thetcbcounter\ \ifx\relax#1\relax\relax\else\textmd{--- \itshape #1}\fi}}
\newtcolorbox{remarque}[1][]{%
  boxbase,colback=remarkcol!6,colframe=remarkcol,
  title={\textsc{Remarque}\ifx\relax#1\relax\relax\else\textmd{ : \itshape #1}\fi}}
\newtcolorbox{attention}[1][]{%
  boxbase,colback=attncol!6,colframe=attncol,
  title={\textsc{Attention}\ifx\relax#1\relax\relax\else\textmd{ : \itshape #1}\fi}}
\newtcolorbox{exemple}[1][]{%
  boxbase,colback=excol!5,colframe=excol,
  title={\textsc{Exemple}\ifx\relax#1\relax\relax\else\textmd{ : \itshape #1}\fi}}
\newtcolorbox{methode}[1][]{%
  boxbase,colback=methcol!7,colframe=methcol,sharp corners,
  title={\textsc{Méthode}\ifx\relax#1\relax\relax\else\textmd{ : \itshape #1}\fi}}
\newtcolorbox{astuce}[1][]{%
  boxbase,colback=primary!7,colframe=primary,
  title={\textsc{Astuce}\ifx\relax#1\relax\relax\else\textmd{ : \itshape #1}\fi}}

\newtcolorbox[auto counter]{exercice}[1][]{%
  boxbase,colback=exercol!4,colframe=exercol,
  title={\textsc{Exercice}~\thetcbcounter\ \ifx\relax#1\relax\relax\else\textmd{--- \itshape #1}\fi}}
\newtcolorbox[auto counter]{correction}[1][]{%
  boxbase,colback=corrcol!4,colframe=corrcol,
  title={\textsc{Corrigé}~\thetcbcounter\ \ifx\relax#1\relax\relax\else\textmd{--- \itshape #1}\fi}}

% ---------- Environnement « où : » ----------
\newenvironment{ou}{%
  \par\smallskip\noindent\textbf{où :}\nopagebreak
  \begin{itemize}[leftmargin=1.8em,topsep=2pt,itemsep=1.5pt,parsep=0pt]}%
  {\end{itemize}}

% ---------- Styles TikZ ----------
\tikzset{
  cat/.style={circle,fill=cationcol,draw=cationcol!50!black,inner sep=0pt,
              minimum size=5mm,font=\bfseries\scriptsize,text=white},
  an/.style={circle,fill=anioncol,draw=anioncol!50!black,inner sep=0pt,
             minimum size=5mm,font=\bfseries\scriptsize,text=white},
  crx/.style={circle,fill=cristalcol,draw=black!50,inner sep=0pt,
              minimum size=5mm,font=\bfseries\scriptsize,text=white},
  ptn/.style={circle,fill=black,inner sep=1.1pt}}

% ---------- En-têtes ----------
\pagestyle{fancy}
\fancyhf{}
\fancyhead[L]{\small\textcolor{primarydark}{\textbf{Fiche 9} — Solubilité et produit de solubilité}}
\fancyhead[R]{\small\textcolor{primarydark}{\textsc{Carnet d'entraînement}}}
\fancyfoot[C]{\small\thepage}
\renewcommand{\headrulewidth}{0.8pt}
\renewcommand{\headrule}{\hbox to\headwidth{\color{primary}\leaders\hrule height \headrulewidth\hfill}}
\usepackage{titlesec}
\titleformat{\section}{\Large\bfseries\color{primarydark}}{\thesection}{0.6em}{}
\titleformat{\subsection}{\large\bfseries\color{primary}}{\thesubsection}{0.6em}{}

% ---------- Bandeau de titre ----------
% #1 = thème/étiquette   #2 = titre principal   #3 = sous-titre
\newcommand{\bandeau}[3]{%
\begin{center}
\begin{tikzpicture}
\node[rectangle,rounded corners=7pt,fill=primary,text=white,
      minimum width=15.6cm,minimum height=1.95cm,align=center,inner sep=8pt]
  {{\large\bfseries #1}\\[3pt]{\Large\bfseries #2}\\[3pt]{\small #3}};
\end{tikzpicture}
\end{center}
\vspace{2mm}}
\begin{document}
\bandeau{Thème 3 — Quotient $Q$ \& précipitation}%
{Corrigé — Niveau Moyen}%
{Réponses détaillées et justifications}

\begin{correction}[magnésium marin — \ce{Mg(OH)2}]
$\ce{Mg(OH)2(s) <=> Mg^{2+} + 2 OH-}$.
\begin{enumerate}[label=\alph*),leftmargin=2em,topsep=2pt,itemsep=2pt]
  \item $Q=[\ce{Mg^2+}][\ce{OH-}]^{2}$.
  \item $Q=(0{,}06)(\num{2{,}0e-3})^{2}=0{,}06\times\num{4{,}0e-6}=\boxed{\num{2{,}4e-7}}$.
        Or $\Kps=\num{9{,}0e-12}$.
  \item $Q=\num{2{,}4e-7}\gg\Kps$ (facteur $\sim\num{2{,}7e4}$) : $Q>\Kps$, il y a
        \textbf{précipitation}, et elle est \textbf{complète} — les ions s'associent en solide jusqu'à
        ce que $Q$ redescende à $\Kps$. \emph{(Réponse « Oui » du \textsc{qcm}\,n\textsuperscript{o}\,9.)}
\end{enumerate}
\end{correction}

\begin{correction}[mélange qui précipite]
\begin{enumerate}[label=\alph*),leftmargin=2em,topsep=2pt,itemsep=2pt]
  \item Volume total $=50+50=\qty{100}{\milli\litre}$, soit un facteur de dilution 2 :
        \[ [\ce{Ag+}]=\dfrac{\num{2{,}0e-4}\times 50}{100}=\qty{1{,}0e-4}{\molar},\quad
           [\ce{Cl-}]=\dfrac{\num{2{,}0e-4}\times 50}{100}=\qty{1{,}0e-4}{\molar}. \]
  \item $Q=[\ce{Ag+}][\ce{Cl-}]=(\num{1{,}0e-4})^{2}=\num{1{,}0e-8}$. Or $\Kps=\num{1{,}8e-10}$, donc
        $Q>\Kps$ (car $\num{1{,}0e-8}=55\times\Kps$).
  \item Oui : $Q>\Kps$, un précipité blanc de \textbf{\ce{AgCl}} se forme.
\end{enumerate}
\end{correction}

\begin{correction}[mélange qui ne précipite pas]
\begin{enumerate}[label=\alph*),leftmargin=2em,topsep=2pt,itemsep=2pt]
  \item Après dilution par 2 : $[\ce{Ag+}]=[\ce{Cl-}]=\qty{5{,}0e-7}{\molar}$, donc
        $Q=(\num{5{,}0e-7})^{2}=\num{2{,}5e-13}$.
  \item $Q=\num{2{,}5e-13}<\Kps=\num{1{,}8e-10}$ : \textbf{aucune précipitation}, tout reste dissous. \\
        \textbf{Comparaison :} même sel, même volumes, mais des concentrations \num{200} fois plus
        faibles font passer $Q$ sous $\Kps$. La précipitation dépend donc \emph{du produit des
        concentrations}, pas seulement de la présence des ions.
\end{enumerate}
\end{correction}

\begin{correction}[seuil de précipitation — \ce{CaF2}]
\begin{enumerate}[label=\alph*),leftmargin=2em,topsep=2pt,itemsep=2pt]
  \item Début de précipitation : $Q=\Kps$, soit $[\ce{Ca^2+}][\ce{F-}]^{2}=\Kps$.
  \item $[\ce{F-}]=\sqrt{\dfrac{\Kps}{[\ce{Ca^2+}]}}=\sqrt{\dfrac{\num{3{,}9e-11}}{\num{1{,}0e-2}}}
        =\sqrt{\num{3{,}9e-9}}\approx\boxed{\qty{6{,}2e-5}{\molar}}$.
        En dessous de cette valeur, $Q<\Kps$ (pas de précipité) ; au-delà, \ce{CaF2} précipite.
\end{enumerate}
\end{correction}

\begin{correction}[sens d'évolution]
\begin{enumerate}[label=\alph*),leftmargin=2em,topsep=2pt,itemsep=2pt]
  \item $Q>\Kps$ : le système évolue dans le sens de la \textbf{formation du solide} (précipitation).
        Les concentrations des ions \textbf{diminuent}, donc $Q$ \textbf{diminue}, jusqu'à ce que
        $Q=\Kps$ (nouvel équilibre saturé).
  \item En diluant fortement, on \emph{diminue} les concentrations, donc $Q$ passe \emph{sous} $\Kps$
        ($Q<\Kps$). Le critère indique alors \textbf{dissolution} : le solide se \textbf{dissout}
        (aucune précipitation). C'est cohérent avec l'idée qu'un sel « insoluble » finit par se
        dissoudre si l'on met assez d'eau.
\end{enumerate}
\end{correction}

\end{document}
